\documentclass{article}
\usepackage[T1]{fontenc}
\usepackage{lmodern}
\usepackage{graphicx}
\usepackage{amsmath,amssymb}
\usepackage[a4paper,margin=2cm]{geometry}
\usepackage[absolute,overlay]{textpos}
\setlength{\TPHorizModule}{1mm}
\setlength{\TPVertModule}{1mm}
\usepackage[indonesian]{babel}
\usepackage{hyperref}
\setlength{\emergencystretch}{2em}
% unit: Halaman 10

\begin{document}

% === Halaman 10 (python/native) ===

10

• Sifat-sifat operasi XOR:

 (i)   a   a = 0 





 (ii)  a  b = b  a 

 (iii) a  (b  c) = (a  b)  c  

  Contoh: 

   (i)  1   1 = 0

   (ii) 1   0 = 0   1 = 1

   (iii) 1  (0  1) = (1  0)  1 = 0

\end{document}
